% LaTeX file for resume

\documentclass[10pt, a4paper]{article}

\usepackage[%top=3.75cm, bottom=3.75cm, left=3.75cm, right=3.75cm
          ]{geometry}
\usepackage{hyperref}
\usepackage{adjustbox}
\usepackage{multicol}
\usepackage{fontspec}
\usepackage{titlesec}
\usepackage{enumitem}
\usepackage{xcolor}
\usepackage{parskip}
\usepackage{microtype}
\usepackage{needspace}

\setlength{\parskip}{3pt}
\hbadness=10000

\hyphenpenalty=10000
\exhyphenpenalty=10000

\setmainfont{Ubuntu}

% Colours
\definecolor{accent}{HTML}{2b2b2b}

% Hyperlink style
\hypersetup{colorlinks=true, urlcolor={blue!60!black}, linkcolor={blue!60!black}}

% Section formatting: small caps, extra space above, no rule
\titleformat{\section}{\large\color{accent}}{}{0em}{}
\titlespacing*{\section}{0pt}{10pt}{6pt}

% Page numbers
\pagestyle{plain}

% Tight lists
\setlist{nosep, leftmargin=0pt, label={}}

\begin{document}

% ---- Header ----
\begin{center}
  {\LARGE\bfseries Thomas Hazell}\\[6pt]
  \url{thomas.hazell@politics.ox.ac.uk} \enspace$\cdot$\enspace \url{https://t-hazell.github.io/}
\end{center}

\vspace{8pt}

% % ---- Research Interests ----
% \section{Research Interests}
% \vspace*{0pt}
% \begin{multicols}{2}
%   Autocracy\\
%   Bureaucracy \& bureaucrats\\
%   Computational social science\\
%   Central Asia\\
%   Procurement \& corruption\\
%   Subnational governance
% \end{multicols}

\section{Education}

\begin{samepage}
  \textbf{DPhil (PhD) in Politics} \hfill 2024 -- \phantom{2027}\\
  \textit{Department of Politics and International Relations, University of Oxford}

Funded by an Economic and Social Research Council Studentship. Supervised by Dr Katerina Tertytchnaya.
\end{samepage}

\begin{samepage}
  \textbf{Visiting Graduate Student} \hfill March -- April 2026\\
  \textit{PSIR, Nazarbayev University, Astana}
\end{samepage}

\begin{samepage}
  \textbf{MPhil in Politics (Comparative Government)} \hfill 2022 -- 2024\\
  \textit{Department of Politics and International Relations, University of Oxford}

Grade: \textit{Distinction}.
\end{samepage}

\begin{samepage}
  \textbf{BA in Philosophy, Politics, and Economics} \hfill 2019 -- 2022\\
  \textit{Lincoln College, University of Oxford}

Grade: \textit{First Class}.
\end{samepage}

\section{Work-in-Progress}

`Personnel appointments in authoritarian bureaucracies'. \textit{Working paper}.

`Social connections and the quality of public procurement'. \textit{In progress}.

`Which merits matter? Political performance, economic meritocracy, and elite loyalty in authoritarian states'. \textit{Working paper}.

`The governance effects of local elections in autocracies'. With Kirill Melnikov and Eleonora Minaeva. \textit{Working paper}.

`Local elections and elite management in authoritarian regimes: evidence from Kazakhstan'. With Eleonora Minaeva and Kirill Melnikov.  \textit{Working paper}.

`Shuffling to co-opt: Subnational governance, co-optation, and political careers in Kazakhstan'. \textit{Working paper}.

\section{Conferences, Workshops \& Talks }

% \textbf{Upcoming}

% `Personnel appointments in authoritarian bureaucracies.' \textit{American Political Science Association Annual Meeting (APSA)}. 5 Septmber 2026. Boston.

% `Elite Selection and Governance in Autocracies: Kazakhstan's Natural Experiment.' (With Kirill Melnikov and Eleonora Minaeva.)  \textit{American Political Science Association Annual Meeting (APSA)}. 4 Septmber 2026. Boston.

% `Personnel appointments in authoritarian bureaucracies.' \textit{Conference of the European Political Science Society (EPSS)}. 18 June 2026. Belfast.

% `Personnel appointments in authoritarian bureaucracies.' \textit{4th APSG Conference on Authoritarian Politics}. 16 June 2026. Authoritarian Political Systems Group, King's College London.

% \textbf{Past}

`Personnel appointments in authoritarian bureaucracies.' 20 March 2026. Department of Political Science and International Relations, Nazarbayev University, Astana.

`Which merits matter? Political performance, economic meritocracy, and elite loyalty in authoritarian states.' \textit{Conference of the European Political Science Association (EPSA)}. 27 June 2025. UC3M, Madrid.

`Personnel appointments in authoritarian bureaucracies.' \textit{3rd APSG Conference on Authoritarian Politics}. 23 June 2025. Authoritarian Political Systems Group, King's College London.

`Personnel appointments in authoritarian bureaucracies.' \textit{3rd Conference ``New Advances in the Political Economy of Development in Eurasia.''} 6 June 2025. Almaty Management University, Almaty.

`Shuffling to co-opt: Subnational governance, patronage, and political careers in Kazakhstan.' \textit{Conference of the Midwestern Political Science Association (MPSA)}. 5 April 2025. Chicago.

`Shuffling to co-opt: Subnational governance, patronage, and political careers in Kazakhstan.' \textit{ESRC Grand Union DTP Annual Conference}. 4 June 2024. University of Oxford, Oxford.

`Shuffling to co-opt: Subnational governance, patronage, and political careers in Kazakhstan.' \textit{Emerging Scholars Workshop on the Study of Challenges in Eurasian Politics and Society}. 24 May 2024. Center for Slavic, Eurasian, and East European Studies, University of North Carolina at Chapel Hill. Presented Online.

\section{Teaching}

\textbf{Undergraduate Tutor} \textit{University of Oxford}\\
Comparative Government, Politics in Russia and the Former Soviet Union

\section{Professional \& Research Experience}

\begin{samepage}
  \textbf{Research Assistant on Non-Violent Repression in Electoral Autocracies}\hfill{}\linebreak September 2024 -- present\\
  \textit{Department of Politics and International Relations, University of Oxford}
\end{samepage}

\begin{samepage}
  \textbf{Election Observer (Short-term), Uzbekistan} \hfill October 2024\\
  \textit{OSCE ODIHR Election Observation Mission}
\end{samepage}

\begin{samepage}
  \textbf{Consultant Researcher (open source methods)} \hfill June -- October 2024\\
  \textit{Ridgeway Information, London}
\end{samepage}

\begin{samepage}
  \textbf{Research Assistant on AUTHLIB (PI: Stephen Whitefield)} \hfill September 2023 -- January 2024\\
  \textit{Department of Politics and International Relations, University of Oxford}
\end{samepage}

\begin{samepage}
  \textbf{Research Assistant to Dr Lenka Buštíková} \hfill May -- September 2023\\
  \textit{Oxford School of Global and Area Studies, University of Oxford}
\end{samepage}

\begin{samepage}
  \textbf{Research Assistant on the Oxford University Economic Recovery Project} \hfill February -- November 2021\\
  \textit{Smith School of Energy and Enterprise, University of Oxford}
\end{samepage}

\begin{samepage}
  \textbf{Intern} \hfill Summer 2021\\
  \textit{Oxford Internet Institute, University of Oxford}
\end{samepage}

\section{Service}

Reviewer for \textit{Journal of Elections, Public Opinion and Parties} and \textit{Communist and Post-Communist Studies}.

\section{Languages}

English (Native). Russian (Intermediate). Qazaq (Beginner).

\section{Computer Skills}

\textbf{R} (Advanced): Modelling \& causal inference methods. Survey data. Visualisation. Managing large datasets. Web scraping/automation. Quarto. Machine learning (\texttt{tidymodels}). GIS (\texttt{sf}).

\textbf{Python} (Basic): Scraping. Data management. Natural language processing/\allowbreak text-as-data/\allowbreak machine learning, including \texttt{torch}. GIS.

\textbf{\textit{Other}}: OpenAI APIs (Batch, Fine Tuning). Git/GitHub. Cloud computing (Azure, Google Cloud). Databases/\allowbreak basic SQL. \LaTeX.

\end{document}
